% New introductory computer-vision exercise; faculty draft.
% Shared layout adapted from the instructor's ME 537 homework sources.
\documentclass[12pt,letterpaper]{article}
\usepackage{mathptmx}
\usepackage[T1]{fontenc}
\usepackage{xcolor,amsmath,amssymb,xurl,listings,enumitem,tikz}
\textwidth=6.5in
\textheight=9in
\oddsidemargin=0in
\evensidemargin=0in
\topmargin=0in
\headheight=0in
\headsep=0in
\parskip=0.08in
\parindent=0in
\setlist[enumerate]{itemsep=0.04in,topsep=0.04in}
\lstset{basicstyle=\ttfamily\footnotesize,breaklines=true,columns=fullflexible,
keepspaces=true,showstringspaces=false,aboveskip=0.04in,belowskip=0.04in}
\newcommand{\R}[2]{{}^{#1}\!R_{#2}}
\newcommand{\T}[2]{{}^{#1}\!T_{#2}}
\newcommand{\courseheader}[2]{%
\setlength{\fboxrule}{2pt}\setlength{\fboxsep}{2mm}%
\fcolorbox{black}{white}{\parbox{\dimexpr\linewidth-2\fboxsep-2\fboxrule\relax}{\centering\textbf{Introduction to Robotics Homework \##1 -- Draft}}}\par
\textbf{#2}\par}
\newcommand{\conventions}{Use right-handed rotations, column vectors, radians,
and meters. $\R{A}{B}$ maps coordinates from frame $B$ into frame $A$;
$\T{A}{B}$ also includes the origin of $B$ expressed in $A$.}
\newcommand{\submission}{\textbf{Submit:} labeled hand derivations and sketches,
your changed Python source files, the requested numerical/visual evidence,
and an AI-use statement (including ``none'' if applicable). Show your reasoning
before using code to check hand calculations. Public checks are intentionally
limited. The instructor supplies the deadline and submission destination.}

\begin{document}
\courseheader{8}{Camera geometry, calibration, and tracking}
\textbf{Schedule:} week 8. Use meters and seconds for world measurements;
image coordinates $(u,v)$ are pixels, with $u$ right and $v$ down.

\textbf{Faculty draft:} The starter code and the object sequence described
below still need to be prepared and reviewed before this can be assigned.

\begin{enumerate}
\item \textbf{Project a point by hand.}
A camera has $f_x=f_y=600$ pixels and principal point $(c_x,c_y)=(320,240)$
pixels. Its optical axes are X right, Y down, and Z forward.
\begin{enumerate}
\item For a point $p_c=[0.12,-0.04,0.80]^T$ m in the camera frame, calculate
its image location using $u=f_x X_c/Z_c+c_x$ and $v=f_y Y_c/Z_c+c_y$.
\item Project $p_c'=[0.24,-0.08,1.60]^T$ m. Explain why one image point
alone cannot tell you a general 3D point's depth.
\end{enumerate}

\item \textbf{Calibrate a planar workspace.}
The two supplied checkerboard images show the same fixed board from two
camera poses. Each square is $0.04$ m wide, and the red dot identifies one
end of the board. The supplied \texttt{corner\_observations.csv} gives
corresponding $(u,v)$ pixels and world-plane $(X,Y)$ meters; four rows are
marked \texttt{fit} and 31 are marked \texttt{check}.
\begin{enumerate}
\item Use the four fit pairs for view 2 with \texttt{cv2.findHomography}
to estimate $H_{P\leftarrow I}$, where
\[
\alpha\begin{bmatrix}X\\Y\\1\end{bmatrix}
=H_{P\leftarrow I}\begin{bmatrix}u\\v\\1\end{bmatrix}.
\]
Explain the input and output units and why normalization by the third
homogeneous coordinate is needed. A CSV loader and point-format helper
will be supplied.
\item Map the 31 \emph{check} corners to the plane. Report the root mean
square and maximum position errors in meters; do not use these corners for
the fit. Show one annotated image identifying fit and check points.
\item Repeat the fit for view 1. Explain why a camera change requires a
new image-to-plane map, even though the board has not moved.
\end{enumerate}
\end{enumerate}

\newpage
\begin{enumerate}[start=3]
\item \textbf{Segment and measure a planar target.}
A supplied fixed-camera sequence will show one red target moving on the
calibration plane. A provided helper will rectify each image onto a grid with
1,000 pixels per meter; its upper-left pixel center represents
$(X,Y)=(0.09,0.29)$ m, with X right and Y up. Use the calibration supplied
for \emph{that sequence's camera}; it may differ from the checkerboard views.
\begin{enumerate}
\item In each rectified BGR image, use \texttt{cv2.inRange} to make a red
mask. The supplied bounds are $(0,0,160)$ through $(90,100,255)$ in BGR
order. Show one frame and its mask. Report the target's bounding box in pixels.
\item Use \texttt{cv2.moments(mask, binaryImage=True)} to calculate the mask
centroid and area. Convert the centroid to world XY using the grid convention,
and convert the area from pixels to square meters. Explain why a circular
target alone does not reveal its in-plane orientation.
\item Return \texttt{None} when the mask has no target pixels. Show one
such frame and explain why $(0,0)$ would be a misleading substitute.
\end{enumerate}

\item \textbf{Estimate motion and discuss limitations.}
\begin{enumerate}
\item By hand, convert grid centroids $(260,260)$ at $t=1.0$ s and
$(300,260)$ at $t=2.0$ s into world positions. Calculate the average XY
velocity in meters per second.
\item Estimate target positions and velocities from the supplied timestamps.
Plot them and compare with the separately supplied evaluation values. Mark
missing detections rather than inventing a position or velocity.
\item Explain why the planar mapping can be inaccurate for a feature above
the calibrated plane or for an image with uncorrected lens distortion.
Why should an extended object's centroid be measured \emph{after} plane
rectification when a metric plane centroid is desired?
\end{enumerate}
\end{enumerate}

\textbf{Project connection:} A calibrated target position and velocity can
help an arm intercept an object on the conveyor. This exercise assumes a fixed
camera and one planar target.

\submission
\end{document}
